\section{Six-player Canadian Fish}
\label{sec:game}

This section specifies the evaluated game variant and then describes the computational and
informational structure that motivates the study. Readers who know the game should still read
\S\ref{sec:game-dialect}, because published descriptions differ on several consequential points and
the rule set studied here is one explicit selection among them. Appendix~\ref{app:dialect} lists
the rule choices against the engine switches that change them.

\subsection{Rules}
\label{sec:game-rules}

Six players form two teams of three, seated alternately so that no two teammates are adjacent. The
deck is 54 cards: the standard 52 plus two jokers. It is partitioned into nine \textbf{half-suits} of
six cards each, being for each suit the low half $2$ through $7$ and the high half $9$ through $A$,
plus a ninth half-suit made of the four eights and the two jokers. Every card belongs to exactly one
half-suit, and each player is dealt nine. The team that takes at least five of the nine half-suits
wins.

Play proceeds by \textbf{asks}. On their turn a player names one opponent and one card. The ask is
legal only if the named card lies in a half-suit that has not yet been claimed, the asker holds at
least one \emph{other} card of that half-suit, the asker does \emph{not} hold the named card, and the
target is an opponent who still holds at least one card. If the target holds the card they must
surrender it and the asker keeps the turn; otherwise the turn passes to the target. The first turn
belongs to the seat after a randomly chosen dealer. When the turn reaches a seat that holds no
cards, that seat names a teammate who still holds cards to receive it; the engine substitutes the
lowest-index such teammate if the named one is ineligible, and if no teammate holds cards the
opposing team enters the forced endgame described in \S\ref{sec:game-dialect}. A seat whose turn it
is and that has no legal ask must declare, and does so on the lowest-index unclaimed half-suit in
which it holds a card. The requirement
that an asker already hold a card of the half-suit makes each legal ask informative about the
asker's holdings, because it turns every ask into a public declaration of partial ownership.

A half-suit is claimed by a \textbf{declaration}. A declaring player names, for each of the six cards
of a half-suit, which member of their \emph{own} team holds it. The claim is all-or-nothing: if every
assignment is correct the declaring team scores the half-suit, and if any assignment is wrong the
opposing team scores it. A declaration is legal only while the half-suit is unclaimed. In the evaluated variant any seat
may declare, in or out of turn and whether or not it holds cards, at any point between two public
events; when several seats propose at the same point, the lowest seat index across both teams
declares first and the others are polled again once that declaration has resolved. A declaration
does not change whose turn it is.

Actions, outcomes, transfers, declarations and hand sizes are public; the identities of the cards
in a player's hand remain private unless a transfer or declaration reveals them. Every ask, its
answer, every card transfer, every declaration and its outcome, and every player's card count are
observed by all six players, so the only hidden information is the initial deal, and the public
record determines every later card position given that deal.

\subsection{Evaluated rule variant}
\label{sec:game-dialect}

Published descriptions of this game document mutually incompatible variants \cite{pagat}, so the rule
set below is one specified variant among the published rule sets. It matches the strategy corpus
this line of work draws on \cite{develin} and the rule set adopted for this implementation. Six of
the eight contested choices Appendix~\ref{app:dialect} lists are engine switches, so their effect
can be measured, and \S\ref{sec:results-dialects} measures the principal comparison over
\vsevenDiaRowsWord{} rule-variant rows: the default and seven variations. The misdeclaration award
is hardcoded and was not varied; the forced-endgame ladder has no flag of its own inside the scored
harness and moves only inside the \texttt{-{}-legacy} bundle, which is why
\S\ref{sec:results-dialects} can estimate it only as a residual. The consequential choices are a
54-card deck in nine half-suits with nine cards per player; declaration legal at any moment, in or out
of turn, by a seat holding cards or not; a misdeclaration awarding the half-suit to the opposing team;
lowest seat index first among simultaneous declarers; an eight-level willingness ladder in the
\textbf{forced endgame} that arises when one whole team is cardless; and a 400-ask action cap whose
residue is adjudicated by physical majority. In the forced endgame the team that still holds cards
must declare every remaining half-suit: the engine polls that team's seats down a ladder of
correctness-probability thresholds, $0.995$, $0.98$, $0.95$, $0.90$, $0.80$, $0.65$ and $0.50$,
followed by a compulsory level at which the first polled seat declares its best guess; at each
level every unclaimed half-suit and every seat of the team is polled in index order, the first
willing seat declares, and the sweep restarts from the top level for the next half-suit. At the
action cap each unclaimed half-suit is awarded to the team physically holding more of its six
cards, and a 3--3 split is awarded to the team holding the lowest-index card of the half-suit; every
half-suit is awarded whole, so a nine-half-suit game cannot tie. In the 48-card, eight-half-suit
variant of \S\ref{sec:results-dialects} each player is dealt eight cards and a 4--4 game is recorded
by the harness as a win for the team seated at index~0; because every deal is played once in each
seating, this is symmetric in expectation across the two configurations, but a tie is not scored as
a half win. No prearranged signalling conventions are available to any policy, and team
communication is limited to the forced-endgame willingness bit and to the successor a cardless
turn-holder names. Appendix~\ref{app:dialect} lists each choice against the switch that changes it.

Two rules are places where the engine is deliberately \emph{more} restrictive than the adopted
rule set allows, and both close coordination channels a policy could otherwise exploit. We evaluated
the two restrictions using the interventions described below. When the turn reaches a cardless seat
whose teammates both hold cards, the rules permit open negotiation over who receives it; the engine
instead has the cardless seat choose unilaterally. A genuine multi-candidate transfer decision
arises \numTransferEvents{} times per game. Replacing the unilateral rule on one team with a
chooser that, given the true deal, hands the turn to the teammate with the most immediately certain
pulls, paired on common deals, changed the mean half-suit score by
$\numTransferDelta \pm \numTransferCI$ half-suits per game (mean and 95\% half-width) over
\numTransferGames{} paired \vfour{} mirror games, in which the two rules diverged in
\numTransferDiverged\% of games. The tested intervention therefore produced no resolved change in
mean half-suit score. This is an oracle-assisted diagnostic rather than a bound on every legal
mechanism: the clairvoyant chooser is a greedy heuristic that does not optimise game value, so the
interval prices the channel for that class of chooser, and it was measured on an earlier policy
rather than on \sestina{}. Similarly, the engine arbitrates simultaneous declarations by lowest seat
index rather than by stated confidence, because a confidence comparison would move a statistic of
one seat's hand to five others. That restriction costs $+\numArbCost{}$~pp of win rate pooled over
\numArbGames{} games on five seeds, and an arbitrator that, given the true deal, selects a correct
declaration from those proposed is worth \numArbOracle{}~pp pooled over three of those seeds. Both
figures measure the effects of two diagnostic arbitration interventions over the fixed proposal
set: the oracle selects a correct proposal when one exists and does not optimise continuation
value, so neither figure bounds the value attainable by arbitration over that set, nor the value of
a policy that changes what it proposes. A half-suit differential and a win-rate edge are different units and are
not converted into one another here.

\subsection{Computational and informational structure}
\label{sec:game-why}

\paragraph{Hidden-state representation.}
Because a card changes hands only through an observed ask, the only thing any player does not know is
the initial deal, and the public record determines every subsequent card position given that deal. In
many hidden-information games the object that must be filtered is the set of reachable
\emph{histories}, for which constructing even one consistent history is FNP-complete in the worst
case, with efficient enumeration available exactly when the public tree is sparse
\cite{solinas-history}. The representation used here reduces hidden-state reconstruction to
constraints on the initial deal: there is exactly one object to reconstruct, and the constraints the
rules impose on it are logical rather than probabilistic. A player told ``I do not hold the seven of
hearts'' knows it, permanently and with certainty. Throughout, a \emph{certificate} is such a hard
constraint on card location derived from a public event.

The research consequence is that Fish permits the constraint-based part of inference to be
\emph{separated} from the policy. Exact enumeration of the deals consistent with the certificates
(\S\ref{sec:inference}) provides a reference for the constraint-based inference computation. It
does not remove uncertainty about the behavioural model used to interpret strategic actions: a
strategically chosen ask can make some feasible deals more likely than others, and that weighting
is a modelling choice the reference calculation does not make (\S\ref{sec:inference-prior}). The
separation isolates, for this paper's negative results, the uncertainty that lies in the policy and
the behavioural model rather than in the certificate arithmetic. It also motivates the belief-model
comparison reported in \S\ref{sec:background-vsix}, in which the exact constraint-based model is a
worse predictor of card location than the deployed approximation with behavioural weighting.

\paragraph{Decentralised team decision-making.}
Three teammates share a payoff, hold different private information, and have no communication channel
beyond their public actions. This is the decentralised control problem the common-information
approach formalises, in which a notional coordinator who knows the common information issues
\emph{prescriptions} mapping each controller's private information to an action
\cite{nayyar2013common}. In Fish the common information is the entire public history, which
induces a common-information belief over the joint private hands. \sestina{} instead maintains an
observer-conditioned belief that combines a seat's private hand with the public history; this is
related to, but not identical with, the coordinator belief used in common-information formulations.
The coordinator's problem is not solvable here, since a prescription maps every possible nine-card hand to
an ask, so the reduction serves as a correctness lens rather than an algorithm. What it does yield
cheaply is a property the cooperative-search literature has to work for \cite{lerer2020sparta}: the
mapping implemented by a \emph{deterministic} procedure reading only public state and its own hand
is common knowledge among the three teammates, with no shared seed and no communication, although
the action it produces at another seat is not, since that seat's hand is private. The determinism
qualifier matters, and \S\ref{sec:agent-tiebreak} is where it becomes an engine switch.

\paragraph{Public policy specification.}
In the extensive-form team-game models this paper cites, a team coordinates through an ex-ante
correlation device: a joint strategy is sampled before play, the opponent may know the distribution
it is sampled from, and the realised sample is private to the team
\cite{celli-gatti,farina-collusion,zhang-tbdag}. A team of three identical, deterministic, published
agents differs from that model in two respects: the policy mapping is public, and there is no
private sampled signal. Knowledge of the policy does not reveal an action determined by a private
hand, but it lets an opponent compute the distribution of a seat's action from the public record
(\S\ref{sec:agent-coordination}). \S\ref{sec:threat} adopts that as the threat model, and
\S\ref{sec:results-adversarial} reports what an adversary that reads the target's policy and
inverts its transcript achieved. Within the search scope described in \S\ref{sec:related-fish}, we
did not identify prior work that addresses this case.

\paragraph{Computational cost.}
A deal resolves in roughly a hundred public events, and on one \vsevenMirrorThreads{}-thread
desktop the engine plays a six-seat mirror of the \emph{blueprint} policy, the non-searching base
policy that the test-time search of \S\ref{sec:agent-search} modifies, at
\vsevenMirrorGamesPerSec{} complete games per second, with the deployed belief running in every
seat; the evaluated configuration, which searches, shows an observed whole-match slowdown of
\vsevenCostFrozenTight$\times$ relative to the non-searching blueprint (\S\ref{sec:results-cost}). This is what makes
the evaluation programme of \S\ref{sec:evaluation} possible: \vsevenBtwoFcheapN{} paired games for a
single headline comparison, two disjoint holdout banks for every claim, a \vsevenPanelMembers{}-member
opponent panel scored identically for four configurations, and \vsevenBfourArmsWord{} separately
fitted adversarial searches all run on one desktop. In a domain where one evaluation cell costs
GPU-hours, the version of this study that reports its own negative results could not have been
afforded. The corresponding liability appears in \S\ref{sec:discussion-external}: low simulation
cost does not provide independent validation.

\paragraph{Planning limitations.}
Three considerations bear on how later results should be read. First, the perfect-information version of
this game offers little to a clairvoyant planner. A player who sees every hand names only cards the
target holds, so never fails an ask, never loses the turn, and declares each half-suit as soon as
their team holds all of it. In \vsevenOmniProbeGames{} mirror games with both teams using the
omniscient probe policy, the non-moving team's score equalled its deal-time locked-half-suit count
in every game (\filepath{research/v06/notes/R8-game-structure.md}). This motivates concern about
weak discrimination by this particular continuation model; the probe does not establish equality or
near-equality of the values of all candidate actions, and no action-level evaluation under
clairvoyant play was recorded. Determinization, the standard recipe for hidden-hand card games
\cite{levy,ginsberg-gib}, would then be expected to reduce here largely to choosing the
ask most likely to hit, blind to the information an ask leaks, the information a refusal supplies,
and the option value of postponing a declaration. \S\ref{sec:related-determinization}
places this against the literature. Second, Fish has no passive continuation: a player must ask or
declare, so there is no equivalent of the check-and-call rollout on which the standard cheap
exploitability probe is built \cite{lisy2017lbr}. \S\ref{sec:discussion-future} names this as the
largest single obstacle between this work and a defensible near-optimality claim. Third, a
deterministic procedure reading only the public state and its own hand can be evaluated by a
teammate on any hypothesised hand; \S\ref{sec:agent-coordination} states the conditions under which
that supports coordination without communication.
